\documentclass[11pt]{article}

\usepackage[utf8]{inputenc}         % Schriftkodierung dieser Datei
\usepackage[english]{babel}          % für deutsche Dokumente

\usepackage{graphicx}               % optional für Grafiken
\usepackage{tabularx}               % optional für Tabellen
\usepackage{multirow}               % optional für Tabellen
\usepackage{url}                % optional für Internet Links

\usepackage[small,bf]{caption2}     % bitte für Bildunterschriften verwenden
\usepackage{parskip}
\usepackage{titlesec}
\usepackage{amsmath}                % optional für Formeln
\usepackage{overpic,subfigure, tikz}

\input{scidefs}
\begin{document}

\title{Low-frequency limit of Wave Field Synthesis}
\date{\today}  % <-- Empty date
\maketitle

\section{High frequency assumptions in Wave Field Synthesis}

Wave Field Synthesis is a well-established sound field reconstruction method allowing the reconstruction of a wave front by driving an array of densely spaced loudspeakers.
The basis of 3D WFS is the Kirchhoff-Helmholtz integral (KHIE) formulation of an arbitrary sound field, expressing the virtual wave front in terms of a double layer potentional surface integral.
In the special geometry of an infinite planar SSD the KHIE can be transformed into a single layer potential containing the driving functions for an infinite planar SSD implicitely, given directly by the gradient of the virtual field.
For arbitrary enclosing SSD surfaces the Kirchhoff approximation can be applied assuming locally planar SSD elements and virtual wave fronts, being a valid assumption in the high frequency region.

Assuming an arbitrary virtual field given by the general polar form
\begin{equation}
    P(\vx,\omega)  = A^P(\vx) \, \te^{ \ti k \Phi^P(\vx)}
\end{equation}
in both cases the involved integral formulation is given by
\begin{equation}
    P(\vx,\omega) \approx \int_S -2 \frac{\partial }{\partial n} P(\vxo,\omega) G(\vx-\vxo,\omega) \, \td S(\vxo) 
\end{equation}
with $\frac{\partial }{\partial n} P(\vxo,\omega)$ denoting the normal derivative of $P(\vx,\omega)$ taken at $\vx = \vxo$, and 
\begin{equation}
    G(\vx-\vxo,\omega) = A^G(\vx) \, \te^{\ti k \Phi^G(\vx)} = \frac{1}{4\pi} \frac{\te^{- \ti k |\vx-\vxo|}}{|\vx-\vxo|}
\end{equation}
being the 3D free field Green's function, where $k = \omega/c$ being the wavenumber.
Again, in case when surface $S$ degenerates to an infinite plane the above expression expresses the target field without any approximation involved, termed as the Rayleigh integral (I),
On the other hand, according to \cite[Eq.(2.7.12)]{Blenstein1975} the approximation holds, when 
\begin{equation}
    k \cdot \rho_{1,2} \gg 1,
\end{equation}
where $\rho_{1,2}$ are the local principal radii of the curved SSD, giving a rough approximation on the frequency limit of validity (and discarding the relative wave front curvature).

For the sake of simplicity the gradient of the target field is usually approximated as
\begin{multline}
    \nabla_\vx  P(\vx,\omega)  = \left( \frac{\nabla_\vx A^P(\vx)}{A^P(\vx)} + \ti k \nabla_\vx \Phi^P(\vx) \right)\, P(\vx,\omega) \approx 
    \ti k \nabla_\vx \Phi^P(\vx) \, P(\vx,\omega)
    = \\ = \ti \vk^P(\vx) \, P(\vx,\omega).
\end{multline}
Where $\vk^P(\vx)$ is the local wavenumber vector, pointing into the direction of the local propagation direction.

Since for a homogeneous medium $|\nabla_\vx \Phi^P(\vx)| = 1$ holds, the approximation is valid when 
\begin{equation}
   \left| \frac{\nabla_\vx A^P(\vx)}{A^P(\vx)}  \right| \ll \frac{\omega}{c}
   \label{eq:HFgradientValidity}
\end{equation}
holds.
The high-frequency Kirchhoff-approximation is, therefore, written as
\begin{equation}
    P(\vx,\omega) \approx -2 \int_S \underbrace{ \ti k_n^P(\vxo) P(\vxo,\omega)}_{D_{3D}(\vxo,\omega)} G(\vx-\vxo,\omega) \, \td S(\vxo) 
    \label{eq:HFRayleigh}
\end{equation}
containing the 3D WFS driving functions $D_{3D}(\vxo,\omega)$ implicitely.


\vspace{2mm}
In practical scenarios the instead of an enclosing surface, or infinite plane, a contour of loudspeakers are employed as SSD elements.
This is termed as 2.5D Wave Field Synthesis.
The backbone of 2.5D WFS is the stationary phase approximation (SPA), stating that for a fixed receiver position $\vx$ the greatest contribution to integral \eqref{eq:HFRayleigh} has one single critical position on the SSD: at which the gradient of the target sound field coincides with the negative gradient of the Green's function, i.e where the phase of the integrand vanishes.
In the present scenario the SPA is used in order to reduce the problem dimensionality from 3D into 2 dimensions, so that the target sound field can be written in the form of a 2D contour integral along $z=0$.
This can be performed if the stationary point is found at $z=0$, i.e. in the plane of interest the target field propagates merely horizontally.
The methodology assumes that the originally 3D surface is invariant in the $z$ dimensions, and the vertical integral at each surface position can be evaluated asymptotically.
Loosely speaking it is assumed that apart from the stationary position (found at $z=0$ for each surface position) the field oscillates rapidly compared to the change of the amplitude, therefore, in the integral rapid oscillations cancel out and the greatest contribution has those parts over the vertical line of integration, where the phase does not change, i.e. its derivative is zero (the stationary position).
This approximation is evidently valid in the high frequency region.

Once the SPA is applied in the vertical dimension (see next section) 2.5D description is found for the virtual sound field, being specific for a fixed receiver position.
WFS driving functions can be already extracted from this formulation.
In practice a further, horizontal SPA assumption is applied to the driving functions: each SSD/contour element dominates the target field along a specific direction, defined by the local wavenumber vector at the SSD element. 
Exploiting this fact one is able control the amplitude of WFS over a predescribed reference curve, which step is out of scope of the present treatise.

As a summary, 2.5D WFS employs the following high frequency assumptions:
\begin{itemize}
    \item Classic formulation apply a high frequency approximation for the required normal derivative of the sound field
    \item 2.5D Kirchhoff approximation applies SPA for dimensionality reduction, being valid in the high frequency region
    \item for the non-planar/linear case the Kirchhoff approximation itself is only valid in the high-frequency region.
\end{itemize}
As a result for violating low frequency requirements the amplitude of the described/synthesized sound field heavily decreases below a given cutoff-frequency, with this frequency limit depending on both the receiver position and the target field model. 
%In the following our aim is to give an analytical approximation on this lower frequency limit.
%In the following first the validity of the SPA is discussed in details, which turns out to give the strictest high frequency requirement from the listed ones.
        
\section{Validity of the Stationary Phase Approximation}

This section discusses the general validity of the SPA in details.
First the derivation of the SPA should be given, up the second non-zero term of the integral's series expansion.
Things get pretty messy here:
\begin{itemize}
\item There's a pretty good derivation given in \cite[p. 114-116]{torresani1995}, yet, seems a bit sloopy.
I made my own derivation based on this one, and shared it with you in spa\_derivation\_torresani.tex.
\item When I asked the chatGPT, it gave me a pretty thorough derivation, however, I did not find it in the Bleistein books in its form.
Maybe I should do more research on it. 
I should also redo this derivation to make sure it did not make any mistakes... Yet, it's a bit complicated. 
You can find it in spa\_derivation\_chatGPT.tex.
\end{itemize}
Anyways, the two derivations give very similar results, aparts from a scaling factor of 24 instead of 2 for the second term of the expansion.
Since the one given by the GPT (with the scaling factor of 2) perfectly gives me the -3dB cutoff frequency in numerical comparisons (see the next section) I tend to belive this one :).
Therefore, according to our virtual friend, the stationary phase approximation up to the second term goes like this:

Assume that we have an oscillatory integral in the general form
\begin{equation}
    I = \int_{-\infty}^{\infty} A(z) \te^{ \ti k \Phi(z) }  \td z.
\end{equation}
It is assumed that there exist a single critical point in the integral path for which
\begin{equation}
    \Phi'(z^*) = 0, \hspace{1cm} \text{and} \hspace{1cm} \Phi''(z^*) \neq 0,
\end{equation}
with $z^*$ termed as the stationary position. 
By expanding both the amplitude and phase into Taylor's series, introducing new variable in order to transform the integral to a Gaussian one, and performing integration the result of integration is given by
\begin{multline}
    \label{eq:result}
    \footnotesize
    I(k)\sim\,  \te^{\ti k \Phi(z^*)}\sqrt{\frac{2\pi}{k|\Phi''(z^*)|}}\,\te^{\ti\frac{\pi}{4}\,\sgn\Phi''(z^*)} \\
    \quad \times \Biggl\{ A(z^*) + \frac{1}{k}\left[\frac{A''(z^*)}{2\Phi''(z^*)}
    -\frac{A(z^*)\Phi^{(4)}(z^*)}{8\left(\Phi''(z^*)\right)^2}+\frac{5A(z^*)\bigl(\Phi^{(3)}(z^*)\bigr)^2}
    {24\left(\Phi''(z^*)\right)^3}\right] + O\Bigl(\lambda^{-2}\Bigr) \Biggr\}.
\end{multline}

\paragraph{Validity of the approximation}
In order to estimate the lower limit for the SPA's validity one may require that the leading term should dominate in the integral approximation, i.e. $|I_0| \gg |I_2|$ should hold.
From \eqref{eq:result}
\begin{equation}
    \scriptsize
    \frac{|I_2|}{|I_0|} =
    \frac{\left| \frac{A''(z^*)}{2\Phi''(z^*)}
    -\frac{A(z^*)\Phi^{(4)}(z^*)}{8\left(\Phi''(z^*)\right)^2}+\frac{5A(z^*)\bigl(\Phi^{(3)}(z^*)\bigr)^2}
    {24\left(\Phi''(z^*)\right)^3}\right|}
    {|A(z^*)|}
    \ll k,
\end{equation}
which can be simplified as
\begin{equation}
    \frac{1}{2}
    \left| \frac{A''(z^*)}{A(z^*) \Phi''(z^*)}
    -\frac{1}{4}\frac{\Phi^{(4)}(z^*)}{\Phi''(z^*)^2}+\frac{5}{12}\frac{\Phi^{(3)}(z^*)^2}
    {\Phi''(z^*)^3}\right|
    \ll k,
            \label{eq:SPAvalidity}
\end{equation}
Furthermore 

\section{Application for the Rayleigh integral}

\subsection{The integral under consideration}
In order to isolate the effect of the SPA from other high frequency assumptions the Rayleigh integral is investigated, by assuming an infinite plane of integration over the $x_0z_0$ plane:
\begin{multline}
    P(\vx,\omega) = - \iint_{-\infty}^{\infty}  2 \frac{\partial}{\partial y} P(\vxo,\omega) \, G(\vx-\vxo,\omega) \, \td x_0 \td z_0 \approx
    \\ \approx \iint_{-\infty}^{\infty}  2 k_y(\vxo) P(\vxo,\omega) \, G(\vx-\vxo,\omega) \, \td x_0 \td z_0,
    \label{eq:Rayleigh}
\end{multline}
or written in the polar form
\begin{multline}
    \scriptsize
    P(\vx,\omega) \approx 2 \iint_{-\infty}^{\infty}  \underbrace{ \ti k_y(\vxo) A^P(\vxo) A^G(\vx-\vxo)}_{A(\vxo)} 
    \te^{-\ti k \overbrace{ \left( \Phi^P(\vxo) + \Phi^G(\vx-\vxo)\right)}^{\Phi(\vxo)} } \, \td x_0 \td z_0,
    \label{eq:RayleighPolar}
\end{multline}
Note that the high-frequency gradient approximation is still applied in order consider only the frequency dependent term in the amplitude, which will help us later to formulate the frequency bound for the SPA validity.
It can be shown that the yielded high frequency limit constitutes a more strict upper frequency bound, than the validity of the gradient approximation given by \eqref{eq:HFgradientValidity}, justifying this decision.

Our aim is to reduce the surface integral along the vertical dimension into a line integral along $x_0$, by performing a stationary phase approximation along $z_0$.
%First the further assumptions in the present application are discussed allowing the significan simpliciation of the validity condition, given by \eqref{eq:SPAvalidity}.

\subsection{Cutoff frequency for the vertical Rayleigh integral:}
As already stated, we restrict our investigation to the horizontal plane $z = 0$ and we assume that in the plane of investigation the target field propagates horizontally, i.e. $k_z^P(x,y,0) = 0$ holds.
As a result the stationary position is identically found at $z_0^* = 0$, independently from the $x_0$ coordinate.
Also, since the field does not propagate into the vertical direction it is ensured that $A'(x_0,y_0,0) = 0$.
Finally, we assume that the source distribution generating the target field is located in the plane of investigation ($z=0$), which results in an even phase function, for which odd derivatives is zero.
These properties also hold for the Green's function.
Therefore, for the amplitude and phase function 
\begin{equation}
    A'(z^*) = 0, \hspace{1cm} \Phi^{(3)}(z^*) = 0
\end{equation}
holds.
With these considerations the validity of SPA \eqref{eq:SPAvalidity} simplifies to
\begin{equation}
    \frac{1}{2}
    \left| \frac{A''(z^*)}{A(z^*) \Phi''(z^*)}
    -\frac{1}{4}\frac{\Phi^{(4)}(z^*)}{\Phi''(z^*)^2}\right|
    \ll k,
            \label{eq:RayleighValidity}
\end{equation}
The first term expresses the fact that in order the oscillations to cancel out it is required that the relative curvature of the amplitude distribution should be significantly smaller, than the curvature of the phase function.
As the cutoff frequency of the approximation in the following we consider the limiting case of equality in \eqref{eq:SPAvalidity} which we define as the -3dB cutoff point of the approximation.

\paragraph{Application for a virtual point source}

For a virtual point source located at $\vxs = [x_s,y_s, 0]^{\mathrm{T}}$ the local wavenumber vector is given by $k_y(x_0,y_0) = k\frac{y_0-y_s}{|\vxo-\vxs|}$ the involved Rayleigh integral reads as
\begin{equation}
    \small
    P(\vx,\omega) \approx -\frac{2 \ti k}{(4 \pi)^2} (y_0-y_s)\iint_{-\infty}^{\infty} \frac{1}{|\vxo-\vxs|^2 |\vx-\vxo|} \te^{-\ti k \left(F |\vxo-\vxs| + |\vx-\vxo| \right)}  \, \td x_0 \td z_0
\end{equation}
where $F = \sgn(y_0-y_s)$ accounts for the description of focused point sources, i.e. sources located in front of the SSD.

For the sake of simplicity in the followings the in-plane distances, and the corresponding curvatures are denoted as
\begin{align}
r_P(x_0) &= \sqrt{(x_0 - x_s)^2 +(y_0 - y_s)^2 }\\  
r_G(x_0) &= \sqrt{(x - x_0)^2 +(y - y_0)^2 } \\
\kappa_P(x_0) &= 1/r_P(x_0) \\
\kappa_G(x_0) &= 1/r_G(x_0) 
\end{align}
The required derivatives in the stationary position ($z_0^* = 0$), without denoting $x_0$-dependency read as
\begin{align}
    A''(x_0,0) &=  - 2\kappa_P^4\kappa_G - \kappa_P^2 \kappa_G^3 \\  
    \Phi''(x_0,0) &=  - \left( F \kappa_P + \kappa_G \right)\\
    \Phi^{(4)}(x_0,0) &= 3 (    F\kappa_P^3 + \kappa_G^3) \\
    \frac{A''(x_0,0)}{A(x_0,0) \, \hat{\Phi}''(x_0,0) } &= - \frac{2 \kappa_P^2 + \kappa_G^2}{F\kappa_P + \kappa_G} \\
    \frac{\Phi^{(3)}(x_0,0)}{|\Phi''(x_0,0)|^2 } &= 3\frac{F \kappa_P^3 + \kappa_G^3}{(F\kappa_P + \kappa_G)^2}
\end{align}
and the lower frequency limit from \eqref{eq:RayleighValidity} assigning a unique cut-off frequency to each SSD element reads as
\begin{equation}
    \omega_0(x_0) = c \left| \frac{2 \kappa_P(x_0)^2 + \kappa_G(x_0)^2}{F\kappa_P(x_0) + \kappa_G(x_0)}
    -
    \frac{3}{4}
    \frac{F \kappa_P(x_0)^3 + \kappa_G(x_0)^3}{(F\kappa_P(x_0) + \kappa_G(x_0))^2}\right| 
        \label{eq:PScutoff}
\end{equation}

\paragraph{Numerical validation}
In order to validate the analytical results given by \eqref{eq:PScutoff} we are able to compare the analytical cutoff frequency by comparing the result of vertical SPA with the actual numerical integration results for each SSD element.
The methodology is given as follows:
\begin{itemize}
    \item The vertical line integral for each SSD element can be evaluted numerically, for the sake of simplicity by using the trapezoid rule, serving as a ground-truth. 
    This results in an integral value for each frequency bin
    \begin{equation}
        I(x_0,\omega) = \sqrt{\omega} \int_{-\infty}^{\infty} \frac{\te^{-\ti \frac{\omega}{c} \left( F k |\vxo-\vxs| + |\vx-\vxo|\right) }}{|\vxo-\vxs|^2 \, |\vx-\vxo|} \td z_0
    \end{equation}
    The term $\sqrt{\omega}$ is left in the integral in order to ensure a flat function above the lower cutoff frequency, merely in the favour of improved vizualization.
    \item The vertical line integral can be also approximated by the SPA, resulting again in an integral value for each frequency bin.
    \item These frequency responses can be compared, and the point where the relative error drops below -3 dB can be estimated numerically. This latter can be compared with the analytical frequency limit given by \eqref{eq:PScutoff}.
\end{itemize}
In the following the integral and its approximation are compared in three different configurations. Furthermore, the estimated and measured cutoff frequencies are compared as the function of $x_0$.
\begin{figure}[]
	\centering
	\begin{overpic}[width = .35\columnwidth ]{figures/setup_1.png}
	\small
	\put(0,0){(a)}
	\end{overpic}    
    \begin{overpic}[width = .35\columnwidth ]{figures/integral_comparison_setup1.png}
        \small
        \put(0,0){(b)}
        \end{overpic}
    \\ \vspace{3mm}
     \begin{overpic}[width = .78\columnwidth ]{figures/cutoff_comparison_setup1.png}
            \small
            \put(0,0){(c)}
            \end{overpic}
	\caption{Rayleigh integral for the sound field of a point source at $\vxs = [0,\, -1,\, 0]^{\mathrm{T}}$, $\vx = [0,\, 1,\, 0]^{\mathrm{T}}$}.
	\label{fig:setup1}
\end{figure}
    
\begin{figure}[]
	\centering
	\begin{overpic}[width = .35\columnwidth ]{figures/setup_2.png}
	\small
	\put(0,0){(a)}
	\end{overpic}    
    \begin{overpic}[width = .35\columnwidth ]{figures/integral_comparison_setup2.png}
        \small
        \put(0,0){(b)}
        \end{overpic}
    \\ \vspace{3mm}
     \begin{overpic}[width = .7\columnwidth ]{figures/cutoff_comparison_setup2.png}
            \small
            \put(0,0){(c)}
            \end{overpic}
            \caption{Rayleigh integral for the sound field of a point source at $\vxs = [-1,\, -0.11,\, 0]^{\mathrm{T}}$, $\vx = [1,\, 0.1,\, 0]^{\mathrm{T}}$}.
            \label{fig:setup2}
\end{figure}
\begin{figure}[]
	\centering
	\begin{overpic}[width = .35\columnwidth ]{figures/setup_3.png}
	\small
	\put(0,0){(a)}
	\end{overpic}    
    \begin{overpic}[width = .35\columnwidth ]{figures/integral_comparison_setup3.png}
        \small
        \put(0,0){(b)}
        \end{overpic}
    \\ \vspace{3mm}
     \begin{overpic}[width = .7\columnwidth ]{figures/cutoff_comparison_setup3.png}
            \small
            \put(0,0){(c)}
            \end{overpic}
            \caption{Rayleigh integral for the sound field of a point source at $\vxs = [0,\, 0.5,\, 0]^{\mathrm{T}}$, $\vx = [0,\, 1,\, 0]^{\mathrm{T}}$}
       \label{fig:setup3}
\end{figure}
Figure \ref{fig:setup1} investigates the most simple case of a point source behind the Rayleigh plane, evaluated in front of the source (same $x$ position) by the Rayleigh integral.
Figure \ref{fig:setup2} shows the more degenerated case of a point source very close to the Rayleigh plane, with the evaluation point being also very close, laterally away from the source.
Finally \ref{fig:setup3} shows the case of a focused source.
In each cases Figure (a) shows the geometry under consideration, Figure (b) shows the vertical integrals at $x_0 = 0$, while Figure (c) depicts the estimated cutoff frequency along the $x_0$ axis.

It is highlighted that for unfocused sources the analytical cutoff frequency \eqref{eq:PScutoff} gives a very good approximation of the cutoff-frequency in each SSD position.
On the other hand for focused sources apart from position in front of the virtual source, large deviations occurs.
This should be investigated that this disceprancy is due to that higher order terms in the SPA also arise significantly in case of focused sources, or I miscalculated something :).
Both explanations may be valid, as when the source gets in front of the SSD the amplitude factor is the very same as it would be if it was located behind the SSD with the same absolute $y_0$ coordinate, but since the phases are subtracted from each other, the phase derivative decrease significantly and SPA does not hold well.
On the other hand I miscalculate things very easily....

Figure \ref{fig:setup3} also dontains some insight to the mechanism of the SPA: Nara asked me in the conference call, what may be the reason behind the apparent ripple in case of focused source WFS.
Apparently, this ripple is alread present in the vertical SPA, so this may be an inherent error of the ray-based approximation. 
Maybe it should worth further discussion.

\section{Application for the Kirchhoff approximation}

\bibliographystyle{unsrt}
\bibliography{dissertation}

\end{document}
